\documentclass[12pt]{article}
\usepackage{amsmath,amssymb,amsthm}
\usepackage[margin=1in]{geometry}

\newtheorem{theorem}{Theorem}
\newtheorem{lemma}[theorem]{Lemma}

\title{Problem 225: Tribonacci Non-divisors}
\author{Project Euler Solution}
\date{}

\begin{document}
\maketitle

\section{Problem Statement}
The Tribonacci sequence is defined by $T_1 = T_2 = T_3 = 1$ and $T_n = T_{n-1} + T_{n-2} + T_{n-3}$ for $n > 3$. Find the 124th odd number that does not divide any Tribonacci number.

\section{Mathematical Foundation}

\begin{theorem}[All Tribonacci Numbers Are Odd]
For all $n \geq 1$, $T_n$ is odd.
\end{theorem}
\begin{proof}
Base cases: $T_1 = T_2 = T_3 = 1$. Inductive step: if $T_{n-3}, T_{n-2}, T_{n-1}$ are odd, then $T_n = T_{n-1} + T_{n-2} + T_{n-3}$ is odd (sum of three odd numbers). By induction, all $T_n$ are odd.
\end{proof}

\begin{theorem}[Pure Periodicity Modulo $m$]
For any positive integer $m$, the sequence $(T_n \bmod m)$ is purely periodic.
\end{theorem}
\begin{proof}
There are $m^3$ possible triples $(T_n, T_{n+1}, T_{n+2}) \bmod m$, so by pigeonhole some triple repeats. The recurrence is invertible ($T_{n-1} = T_{n+2} - T_{n+1} - T_n$), so backward uniqueness implies pure periodicity.
\end{proof}

\begin{lemma}[Period Bound]
The period of $(T_n \bmod m)$ is at most $m^3$.
\end{lemma}
\begin{proof}
There are $m^3$ possible triples, so repetition occurs within $m^3 + 1$ steps.
\end{proof}

\begin{theorem}[Divisibility Criterion]
An odd number $m$ divides some $T_n$ if and only if $0$ appears in one complete period of $(T_n \bmod m)$.
\end{theorem}
\begin{proof}
By pure periodicity, $0$ appears in the sequence iff it appears in a single period.
\end{proof}

\begin{lemma}[Cycle Detection]
Starting from $(1,1,1)$, compute successive triples modulo $m$. If $T_n \equiv 0$ is encountered, $m$ is a divisor. If $(1,1,1)$ recurs without $0$, $m$ is a non-divisor.
\end{lemma}
\begin{proof}
The sequence is purely periodic with initial triple $(1,1,1)$, so the first recurrence of $(1,1,1)$ marks one complete period.
\end{proof}

\section{Algorithm}
For each odd $m = 1, 3, 5, \ldots$:
\begin{enumerate}
\item Compute $T_n \bmod m$ starting from $(1,1,1)$.
\item If $0$ appears, skip. If $(1,1,1)$ recurs, increment counter.
\item Stop when counter reaches 124.
\end{enumerate}

\section{Complexity Analysis}
\begin{itemize}
\item \textbf{Time:} $O(m^3)$ worst-case per candidate, but $O(m)$ in practice. Total: $O(m_{\max}^2)$ empirically.
\item \textbf{Space:} $O(1)$ per candidate.
\end{itemize}

\section{Answer}

\[
\boxed{2009}
\]

\end{document}
